\chapter{Data}
\label{ch:data}

\section{Source and pre-processing}

The data set used in this dissertation is the Berkeley Earth Global Land Temperatures by City dataset, which is publicly available via Kaggle. Only a single monthly snapshot $\texttt{2003-07-01}$ (July 2003) is used in this dissertation, so that the focus is on the spatial reconstruction problem, which is similar to the previous dissertation by Will Thomas. Berkeley Earth city-level temperatures are not raw thermometer measurements. Instead, the original station data are processed using spatial interpolation, homogenisation, and other processing methods. Therefore, each city temperature is a processed estimate when interpreting the residuals in the results section (\cref{ch:results}). Residuals here measure the difference between the spherical harmonic reconstruction and the Berkeley Earth processed temperature field. So, they should not be interpreted as the error relative to a completely raw physical measurement, which is discussed later as a limitation in the discussion chapter (\cref{ch:discussion}).

The large raw data set is processed in 50,000-row chunks, keeping only observations corresponding to July 2003 to reduce memory usage. The latitude, longitude, average temperature, city, and country variables were retained, while rows with missing average temperature were removed. Latitude and longitude strings like \texttt{"34.56N"} or \texttt{"12.3W"} are converted to colatitude $\theta$ and longitude $\phi$ in radians. For northern latitudes, $\theta = (90^\circ - \mathrm{lat})\cdot\pi/180$, while the complementary form is used for southern latitudes. The longitude $\phi$ is wrapped into $[0,2\pi)$ eastward. After cleaning, $n=3,510$ cities remain, with each observation now containing a single July average temperature (in $^\circ\mathrm{C}$), a colatitude, and a longitude. A summary of the resulting dataset is provided in \Cref{tab:data-summary}. 

\begin{table}[htbp]
\centering
\caption{Summary statistics of the cleaned dataset ($n=3{,}510$ cities, July 2003).}
\label{tab:data-summary}
\begin{tabularx}{\textwidth}{l>{\raggedright\arraybackslash}X}
\toprule
Quantity & Value \\
\midrule
Number of cities & 3,510\\
Observed temperature range & $-0.75^\circ\mathrm{C}$ to $38.38^\circ\mathrm{C}$ \\
Latitude range covered & approx.\ $55^\circ\mathrm{S}$ to $70^\circ\mathrm{N}$ \\
Poleward gaps & no observations south of $\sim\!55^\circ\mathrm{S}$ or north of $\sim\!70^\circ\mathrm{N}$ \\
Ocean coverage & none (land-based city stations only) \\
Train / Validation / Calibration / Test split & 2,106 / 702 / 351 / 351\\
\bottomrule
\end{tabularx}
\end{table}
These observations are not uniformly distributed across the sphere. Some regions have relatively dense city coverage, such as Europe, Eastern North America, and Eastern China, while other regions have sparser coverage, such as parts of South America, Africa, and Central Asia. There are almost no observations over the oceans, and limited coverage at high latitudes, approximately beyond $70^\circ\mathrm{N}$. This uneven spatial distribution is important for the dissertation because the proposed method estimates local prediction difficulty by using spatial neighbourhoods (\cref{ch:methodology}). 
 
There are two things to be careful about when interpreting the uncertainty maps (\cref{ch:results}). First, this dataset contains almost no observations over the oceans. This means that the predictions over large oceanic regions are mainly extrapolations from land-based observations, rather than interpolations between nearby observations. Since the difficulty estimator is based on residuals from nearby training observations and there are no nearby observations to estimate local difficulty reliably, it may not fully capture the difficulty of these ocean predictions. This is different from the Andes example, where observations are present but are relatively sparse. Therefore, local difficulty can still be estimated from nearby observations. This distinction is important when interpreting the uncertainty maps (\cref{sec:geography}). Second, the dataset is a collection of cities, rather than a random sample of points on the sphere. Cities themselves are not randomly distributed, and their locations can be related to physical and geographical factors such as elevation, terrain, population distribution, and accessibility. Therefore, the spatial distribution of observations may also affect the uncertainty results. This should be considered when interpreting the geographic concentration of high-uncertainty locations in the Andes.
 
\section{Data splitting}

The data have to be split to ensure that the different stages of the conformal prediction procedure do not use the same observations. In particular, the calibration data are not used to fit the predictor, estimate $\hat{\sigma}(\cdot)$, or select hyperparameters. Therefore, the full dataset is divided into four separate subsets:

\begin{itemize}
\item \textbf{Training} ($60\%$, $n=2{,}106$): used to fit the spherical harmonic regressor, estimate the difficulty function $\hat{\sigma}(\cdot)$, and perform 5-fold cross-validation. In each fold, the model is fitted on four folds and used to make predictions on the held-out fold. This produces out-of-fold residuals that are used to estimate local prediction difficulty. This avoids using overly optimistic in-sample residuals.

\item \textbf{Validation} ($20\%$, $n=702$): used for hyperparameter selection in the degree-robustness experiment. For each tested spherical harmonic degree $L$, the validation set is used to select $\lambda$ according to the two selection criteria (RMSE criterion and conformal-efficiency criterion). Validation data are separate from the training, calibration, and test sets. 

\item \textbf{Calibration} ($10\%$, $n=351$): used to calculate conformal nonconformity scores and the corresponding conformal quantile $\hat{q}_{\alpha}$ after the base predictor and difficulty estimator have been fixed.

\item \textbf{Test} ($10\%$, $n=351$): used only for the final evaluation of prediction accuracy and coverage analysis.

\end{itemize}
 
This four-way split separates the different stages of the proposed method and avoids data leakage. This separation is important for maintaining the validity of the split-conformal procedure. In particular, the calibration scores should be calculated using a predictor that has not been fitted using the calibration data. If the calibration data were used to choose $\lambda$, $k$, or $L$, the fitted predictor would depend on the same observations that are subsequently used to calculate the conformal scores. The conformal scores would no longer be independent of model selection, and standard split-conformal validity would not directly apply.

The training data are used to obtain the difficulty estimator $\hat{\sigma}(x)$, while the calibration data are kept separate from this step. As a result, $\hat{\sigma}(x)$ is fixed before calculating the calibration scores. The calibration scores are then calculated as
\begin{equation}
S_i =
\frac{|y_i-\hat{y}(x_i)|}
{\hat{\sigma}(x_i)},
\qquad i=1,\ldots,n_{\mathrm{cal}},
\end{equation}

where $x_i$ denotes the geographical location of the $i$th calibration observation and $y_i$ denotes its observed temperature. An empirical quantile is calculated from these scores which is then used to construct prediction intervals for the test observations and other prediction locations.



% ============================================================
